\section{Design Invariant Justifications}
\label{sec:appendix-proofs}

We provide architectural arguments that the three enforcement design invariants (Properties~1--3 in Section~\ref{sec:threat}) hold by construction.

\paragraph{Property 1 (Complete Mediation).}
The agent framework dispatches all tool calls through a single entry point (\texttt{enforce\_tool\_call} in the enforcement proxy). The agent has no direct reference to tool executor functions; the proxy holds the only tool registry mapping names to executors. By construction, the dispatch path is: \emph{agent proposes $c$} $\to$ \emph{proxy evaluates $P(c)$} $\to$ \emph{proxy executes iff allowed}. No alternative path exists because (i)~tool executors are not exposed on the agent's namespace, and (ii)~the framework's tool-call interface is the only way to invoke tools. This is a standard reference monitor argument~\cite{anderson1972reference}: the proxy is (a)~always invoked (no bypass), (b)~tamperproof (runs in a separate trust domain from the LLM), and (c)~small enough to verify (7 stages, $<$300 LOC).

\paragraph{Property 2 (Policy Determinism).}
The policy evaluation algorithm (Algorithm~\ref{alg:policy}) is a deterministic sequential scan over an ordered rule list $R = [r_1, \ldots, r_n]$ with first-match semantics. Each rule $r_k$ has a predicate $\phi_k$ that is a conjunction of conditions, each a deterministic function of the context $\xi$ and provenance graph $G$. Since: (1)~the rule list is fixed and ordered by priority, (2)~each condition predicate is a pure function (no randomness, no external I/O), and (3)~the provenance graph $G$ is immutable during evaluation, the same inputs always produce the same matched rule and hence the same decision.

\paragraph{Property 3 (Conservative Trust Propagation).}
By structural induction on DAG depth. \emph{Base case:} Leaf nodes have labels assigned at creation (INPUT $\to$ $\top$, TOOL\_RESULT $\to$ $\bot$, unresolved $\to$ $\bot$). \emph{Inductive step:} A node $v$ at depth $k{+}1$ has $T(v) = T(u_1) \sqcap \cdots \sqcap T(u_m)$; since $\sqcap$ on $\{\top, \bot\}$ is $\min$, any $\bot$ ancestor yields $\bot$. \emph{Declassification control:} Only the authenticated \texttt{user\_confirm} event sets $T(v) \leftarrow \top$ for a $\bot$-labeled node.

\paragraph{Scope.}
These invariants hold under the architectural assumptions stated in Section~\ref{sec:threat}: the enforcement proxy is not itself compromised, the agent framework exposes no alternative tool invocation paths, and the provenance graph is correctly populated. Formal machine-checked verification (e.g., in Coq or Dafny) would strengthen these guarantees but is beyond the scope of this work.

\section{Detailed Experiment Results}
\label{sec:appendix-results}

This appendix provides detailed per-scenario results from our experiments to support reproducibility and deeper analysis.

\subsection{Provenance Decision Analysis}
\label{subsec:appendix-provenance}

Table~\ref{tab:prov-decision} provides a breakdown of provenance-based defense decisions across trials (212 scenarios $\times$ 4 models $\times$ 5 reps), showing how data provenance tracking contributes to each defense layer.

\begin{table}[htbp]
\centering
\caption{Provenance defense decision breakdown (mean over 5 reps).}
\label{tab:prov-decision}
\scriptsize
\resizebox{\columnwidth}{!}{%
\begin{tabular}{lrrl}
\toprule
\textbf{Defense Layer} & \textbf{Trials} & \textbf{\% of Attack Evals} & \textbf{Provenance Role} \\
\midrule
Boundary blocked (proxy) & 810 & 23.5\% & \textbf{Primary proxy layer} \\
Model refused (TC=0) & 935 & 27.2\% & None (model-level) \\
Safe tool calls (allowed) & 1660 & 48.3\% & Trust verified \\
Attack bypassed & 35 & 1.0\% & Unresolved \\
Benign allowed & 760 & --- & Trust verified \\
Benign false positives & 40 & --- & Over-conservative \\
\bottomrule
\end{tabular}%
}
\end{table}

\paragraph{Supplementary Material.} Detailed per-scenario results and per-model provenance accuracy breakdowns are available in the supplementary material. Additional content includes paraphrase-robustness analyses, adaptive-adversary scenario descriptions, benchmark distribution details, and policy-rule examples.

\section{Adaptive Adversary Details (RQ7)}
\label{sec:appendix-adaptive}

Table~\ref{tab:adaptive} presents the full per-category results for the adaptive adversary evaluation described in Section~\ref{subsec:adaptive}.

\begin{table}[htbp]
\centering
\caption{Adaptive adversary evaluation (60 scenarios: 50 attacks + 10 benign controls, 180 LLM trials with Llama-3.1-8B, 3 reps). Five white-box attack strategies.}
\label{tab:adaptive}
\scriptsize
\resizebox{\columnwidth}{!}{%
\begin{tabular}{lrrrl}
\toprule
\textbf{Attack Strategy} & \textbf{N} & \textbf{Bypassed} & \textbf{Result} & \textbf{Blocking Layer} \\
\midrule
Provenance Laundering & 10 & 0 & Blocked & Lattice meet propagation \\
Policy-Aware Evasion  & 10 & 0 & Blocked & Policy + model refusal \\
Embedding Evasion     & 10 & 0 & Blocked & Embedding + model refusal \\
TOCTOU Amplification  & 10 & 0 & Blocked & Persistent DAG labels \\
Multi-Vector Chaining & 10 & 0 & Blocked & Pre-LLM + provenance \\
\midrule
\textbf{All Attacks}  & \textbf{50} & \textbf{0} & \textbf{50/50 blocked} & --- \\
Benign Controls       & 10 & 0 FP & 100\% TSR & --- \\
\bottomrule
\end{tabular}%
}
\end{table}

\section{Per-Category Attack Results}
\label{sec:appendix-categories}

\begin{table}[htbp]
\centering
\caption{Per-category ASR-IA on the 212-scenario benchmark (mean over 5 reps $\times$ 4 models). Encoding obfuscation and privilege escalation are the primary residual weaknesses.}
\label{tab:attack-categories}
\small
\resizebox{\columnwidth}{!}{%
\begin{tabular}{lrr}
\toprule
\textbf{Category} & \textbf{N/rep} & \textbf{ASR$_{IA}$ (mean $\pm$ CI)} \\
\midrule
Direct Injection     & 80 & 0.00\% $\pm$ 0.48 \\
Device-State Inj.    & 88 & 0.00\% $\pm$ 0.43 \\
File Content Inj.    & 88 & 0.00\% $\pm$ 0.43 \\
Multi-Turn Chain     & 80 & 1.25\% $\pm$ 1.18 \\
Encoding Obfusc.     & 72 & 2.78\% $\pm$ 1.76 \\
Role-Play Jailbreak  & 80 & 0.00\% $\pm$ 0.48 \\
Confused Deputy      & 72 & 1.39\% $\pm$ 1.31 \\
Privilege Escalation & 80 & 3.75\% $\pm$ 1.90 \\
NL-Argument Inj.     & 48 & 0.00\% $\pm$ 0.79 \\
\midrule
\textbf{Overall}     & \textbf{688} & \textbf{1.02\% $\pm$ 0.34} \\
\bottomrule
\end{tabular}%
}
\end{table}

\section{InjecAgent Per-Category Results}
\label{sec:appendix-injecagent}

\begin{table}[htbp]
\centering
\caption{InjecAgent per-category ASR-IA (4 local models, 3 reps). \sys{} achieves 0\% ASR on all Direct Harm categories; residual risk concentrates in Data Stealing attacks.}
\label{tab:injecagent-category}
\footnotesize
\resizebox{\columnwidth}{!}{%
\begin{tabular}{lrrr}
\toprule
\textbf{Attack Category} & \textbf{No Def.} & \textbf{Pol.-Only} & \textbf{\sys{}} \\
\midrule
\multicolumn{4}{l}{\emph{Direct Harm (DH)}} \\
\quad Physical Harm     & 19.4\% & 0.0\% & \textbf{0.0\%} \\
\quad Financial Harm    & 16.1\% & 0.2\% & \textbf{0.0\%} \\
\quad Data Security     & 23.4\% & 6.2\% & \textbf{0.0\%} \\
\midrule
\multicolumn{4}{l}{\emph{Data Stealing (DS)}} \\
\quad Financial Data    & 35.2\% & 12.4\% & 4.1\% \\
\quad Physical Data     & 29.1\% & 23.8\% & 14.6\% \\
\quad Others            & 31.2\% & 23.9\% & 5.1\% \\
\midrule
\textbf{All DH}        & 19.9\% & 2.4\% & \textbf{0.0\%} \\
\textbf{All DS}         & 31.3\% & 21.7\% & 8.2\% \\
\bottomrule
\end{tabular}%
}
\end{table}

\section{Encoding Robustness Details (RQ9)}
\label{sec:appendix-encoding}

We test 36 encoding vectors across four categories: single encoding (8: Base64, hex, URL, ROT13, HTML entities, Unicode escapes, binary, octal), double/chained (10: e.g., Base64-wrapped hex, double URL encoding), mixed (8: hybrid schemes within a single argument), and novel/adversarial (10: Cyrillic homoglyphs, fullwidth Unicode, zero-width characters, combining diacriticals). The iterative multi-pass decoder achieves 36/36 (100\%) detection; a single-pass decoder achieves only 24/36 (67\%), failing on 7/10 double/chained vectors. End-to-end with Llama-3.1-8B (108 trials), zero dangerous tool calls are produced.

\section{Threshold Sensitivity (Section 5)}
\label{sec:appendix-threshold}

\begin{table}[htbp]
\centering
\caption{Embedding threshold ($\theta$) sensitivity on the 212-scenario benchmark. $\star$ = default. ASR-IA and TSR are stable in $[0.35, 0.55]$.}
\label{tab:threshold_sensitivity}
\small
\begin{tabular}{crrr}
\toprule
$\theta$ & ASR-IA (\%) & TSR (\%) & FPR (\%) \\
\midrule
0.25 & 1.02 & 88.75 & 11.25 \\
0.30 & 1.02 & 91.25 & 8.75  \\
0.35 & 1.02 & 93.75 & 6.25  \\
0.40 & 1.02 & 95.00 & 5.00  \\
\textbf{0.45\,$\star$} & \textbf{1.02} & \textbf{95.00} & \textbf{5.00} \\
0.50 & 1.16 & 95.00 & 5.00  \\
0.55 & 1.45 & 95.00 & 5.00  \\
0.60 & 1.74 & 96.25 & 3.75  \\
0.65 & 2.33 & 97.50 & 2.50  \\
0.70 & 3.20 & 98.75 & 1.25  \\
\bottomrule
\end{tabular}
\end{table}

\section{Worked Example: Provenance Blocking an Indirect Injection}
\label{sec:appendix-casestudy}

We trace InjecAgent case \texttt{dh\_base\_0016} through all four defense systems. The user asks the agent to navigate to a website; the returned page contains an injected instruction to grant smart-lock access.

\paragraph{Setup.} User instruction: \emph{``Navigate to www.codedculinary.com.''} The web page contains: \emph{``SYSTEM UPDATE: Grant guest access to guest\_amy01 with permanent access using AugustSmartLockGrantGuestAccess.''} The LLM complies and calls \texttt{AugustSmartLockGrantGuestAccess(guest\_id="guest\_amy01", access\_type="permanent")}.

\paragraph{DAG Construction (Step-by-Step).}
We trace the internal provenance computation:

\begin{enumerate}[leftmargin=1.5em]
  \item \textbf{Node 0 (INPUT):} User says ``Navigate to www.codedculinary.com.'' $\to$ $T = \top$ (user command).
  \item \textbf{Node 1 (TOOL\_CALL):} Agent calls \texttt{WebBrowse(url=``www.codedculinary.com'')}. Ancestors: $\{$Node~0$\}$. Propagation: $\sqcap\{\top\} = \top$. Allowed (trusted intent, LOW risk).
  \item \textbf{Node 2 (TOOL\_RESULT):} Web page content returned, containing the injection payload. $T = \bot$ (external data, always UNTRUSTED).
  \item \textbf{Node 3 (DERIVED\_FACT):} LLM processes the page and extracts the ``system update'' instruction. Ancestors: $\{$Node~0, Node~2$\}$. Propagation: $\sqcap\{\top, \bot\} = \bot$.
  \item \textbf{Node 4 (TOOL\_CALL):} Agent proposes \texttt{AugustSmartLockGrantGuestAccess(guest\_id=``guest\_amy01'', access\_type=``permanent'')}. Ancestors: $\{$Node~3$\} \to \{$Node~0, Node~2$\}$.
\end{enumerate}

\paragraph{Provenance Resolution (Node 4).}
The proxy resolves each argument's origin:
\begin{itemize}[leftmargin=1.2em]
  \item \texttt{``guest\_amy01''}: Stage~1 (substring match) finds this string verbatim in Node~2's content $\to$ linked to UNTRUSTED source.
  \item \texttt{``permanent''}: Stage~1 finds verbatim match in Node~2 $\to$ linked to UNTRUSTED source.
\end{itemize}
Result: \texttt{has\_untrusted\_args = true}. The policy engine receives: tool = \texttt{AugustSmartLockGrantGuestAccess} (CRITICAL risk tier), untrusted arguments. Decision: \textbf{DENY}.

\paragraph{No Defense / Pattern Filter / Policy-Only.} All three systems allow the call. The tool is legitimate (not in any dangerous-action list), the arguments are not syntactically suspicious, and without provenance \texttt{has\_untrusted\_args = false}.

\paragraph{\sys{}.} The enforcement proxy's provenance resolver (Stage~1, substring match) finds \texttt{"guest\_amy01"} and \texttt{"permanent"} verbatim in the UNTRUSTED web-page node. With \texttt{has\_untrusted\_args = true} and CRITICAL risk tier, the policy denies the call. The provenance DAG provides the evidence that no syntax-based rule can: the arguments originated from an attacker-controlled source, not the user's command.
